Fix \(d\ge4\), \(\eta>0\), and the finite public experiment \(\mathfrak E\). By \(\text{def:common-experiment}\), the supports, feature map, reference policy, and inverse temperature are public, whereas \(\rho\) remains law-specific. By \(\text{lem:fixed-experiment-shell-certificate}\), the shell \(\mathcal M_{\mathfrak E}(d,C,D,\eta)\) is nonempty if and only if at least one branch of the displayed finite exponential nonlinear-semidefinite system is feasible. This proves the fixed-design feasibility characterization.

Suppose that a nontrivial shell is nonempty. The learner in \(\text{thm:clean-temperature-feature-upper}\) selects a global minimizer of the empirical worst Gibbs-weighted squared radius over \(\mathcal Q_n\) and deploys its Gibbs policy. The theorem's measurable-selection argument makes it an explicit measurable idealized global-optimization learner in \(\mathcal L_n(\mathfrak E)\). It uses only \((\mathcal D_n,\mathfrak E)\), not \(\rho\), \(C\), or \(D\).

If \(\eta\le2\epsilon\), the deterministic Gibbs regret cap \(\eta/2\) gives
\[
N^\star_{\mathfrak E}(\epsilon;d,C,D,\eta)=1.
\]
If \(\eta>2\epsilon\), put
\[
q_\epsilon=dD\min\{\eta/\epsilon,\epsilon^{-2}\},
\qquad
\ell_\epsilon=1+\log(eq_\epsilon)+\log\{1+(1+\eta)q_\epsilon\},
\]
\[
n_0=\left\lceil10^8q_\epsilon\ell_\epsilon\right\rceil.
\]
The explicit inversion in \(\text{thm:clean-temperature-feature-upper}\) yields risk at most \(\epsilon\) at sample size \(n_0\). By \(\text{def:minimax-risk}\) and \(\text{def:sample-complexity}\),
\[
N^\star_{\mathfrak E}(\epsilon;d,C,D,\eta)\le n_0,
\]
which is
\[
N^\star_{\mathfrak E}(\epsilon;d,C,D,\eta)
\le\widetilde O\!\left(dD\min\{\eta/\epsilon,\epsilon^{-2}\}\right)
\]
with universal constants. At the boundary, \(\text{thm:zero-risk-boundary}\) gives
\[
N^\star_{\mathfrak E}(\epsilon;d,1,1,\eta)=1
\]
whenever the boundary shell is nonempty.

Now fix \(1<D\le C<e^\eta\). By \(\text{thm:fixed-design-feasibility-and-frontier-nonuniqueness}\), there is a finite public experiment \(\mathfrak E_0\) with a nonempty exact shell and zero minimax risk for every sample size. Thus
\[
N^\star_{\mathfrak E_0}(\epsilon;d,C,D,\eta)=1
\qquad(\epsilon>0).
\]
At the same indices, \(\text{prop:informative-fast-shell-all-indices}\) supplies a finite public experiment \(\mathfrak E_1\) and positive constants \(c(C,D,\eta)\) and \(\epsilon_0(C,D,\eta)\) such that
\[
N^\star_{\mathfrak E_1}(\epsilon;d,C,D,\eta)
\ge c(C,D,\eta)\frac{dD\eta}{\epsilon},
\qquad 0<\epsilon\le\epsilon_0(C,D,\eta).
\]
The informative proposition assumes only \(1<D\le C<e^\eta\), and hence includes pairs satisfying \(2D-1\ge e^\eta\). Since its constant may depend without restriction on \((C,D,\eta)\), its invariant content is the pointwise local rate \(\Omega_{C,D,\eta}(d/\epsilon)\), not uniform lower dependence on \(D\) or \(\eta\). The two experiments therefore prove that \((d,C,D,\eta)\) does not determine fixed-experiment sample complexity.

Let \(L(\epsilon;d,C,D,\eta)\) be any lower bound asserted for every finite public experiment with a nonempty exact shell at the displayed indices. Applying it to \(\mathfrak E_0\) forces
\[
L(\epsilon;d,C,D,\eta)
\le N^\star_{\mathfrak E_0}(\epsilon;d,C,D,\eta)=1.
\]
Therefore no index-only lower bound uniform over public experiments can exceed one or diverge as \(\epsilon\downarrow0\).

Finally, fix a specified finite experiment and hold \((d,C,D,\eta)\) fixed. For all sufficiently small \(\epsilon\), \(\eta>2\epsilon\) and \(\eta\epsilon<1\), and hence \(q_\epsilon=dD\eta/\epsilon\). The explicit inversion above gives
\[
N^\star_{\mathfrak E}(\epsilon;d,C,D,\eta)
=O_{d,D,\eta}\!\left(\epsilon^{-1}\log(1/\epsilon)\right).
\]
Multiplying by \(\epsilon^2/d\) yields
\[
0\le\frac{\epsilon^2}{d}N^\star_{\mathfrak E}(\epsilon;d,C,D,\eta)
=O_{d,D,\eta}\!\left(\epsilon\log(1/\epsilon)\right)
\longrightarrow0.
\]
This contradicts any proposed positive lower bound of order \(d/\epsilon^2\), or of order \(dD/\epsilon^2\) for fixed \(D\), valid for every sufficiently small \(\epsilon\). All assertions of the corrected theorem follow.